% ============================================================
%  main.tex  —  Overleaf master file (wmt-template compliant)
%  Paper: Multi-Objective Optimization of 2E-MD-MD-MP-CVRP-SPD
%         using Hybrid NSGA-III (HNSGA-III)
% ============================================================
\documentclass[11pt, a4paper]{article}

% ====== Preamble ======
\usepackage{wmt-template}
\usepackage[section]{placeins}   % keeps floats within their section

\addbibresource{references.bib}
\graphicspath{{./}{./figures/}}

\algnewcommand{\To}{\textbf{to }}
\algnewcommand{\Step}{\textbf{step }}

\studentname{[Author Name]}

% ── Custom math shorthands ────────────────────────────────────────
\newcommand{\KFE}{K_{\mathrm{FE}}}
\newcommand{\KSE}{K_{\mathrm{SE}}}
\newcommand{\NFE}{N_{\mathrm{FE}}}
\newcommand{\NSE}{N_{\mathrm{SE}}}
\newcommand{\QFE}{Q_{\mathrm{FE}}}
\newcommand{\QSE}{Q_{\mathrm{SE}}}
\newcommand{\cFE}{c_{\mathrm{FE}}}
\newcommand{\cSE}{c_{\mathrm{SE}}}
\newcommand{\eSE}{e_{\mathrm{SE}}}
\newcommand{\vSE}{v_{\mathrm{SE}}}
\newcommand{\vFE}{v_{\mathrm{FE}}}
\newcommand{\CFE}{C_{\mathrm{FE}}}
\newcommand{\CSE}{C_{\mathrm{SE}}}
\newcommand{\EFE}{E_{\mathrm{FE}}}
\newcommand{\ESE}{E_{\mathrm{SE}}}

\algrenewcommand\algorithmicindent{1.2em}
\numberwithin{equation}{section}

\title{\textbf{Multi-Objective Optimization of the Two-Echelon Multi-Depot\\
Multi-Day Multi-Period Capacitated Vehicle Routing Problem\\
with Simultaneous Pickup and Delivery Using Hybrid NSGA-III}}

\author{[Author Names]}
\date{}

% ============================================================
\begin{document}
\maketitle
\thispagestyle{fancy}

% ─────────────────────────────────────────────────────────────
\begin{abstract}
This study addresses the two-echelon multi-depot multi-day multi-period
capacitated vehicle routing problem with simultaneous pickup and delivery
(2E-MD-MD-MP-CVRP-SPD), a challenging logistics optimization problem that
simultaneously minimizes total transportation cost, vehicle count, and carbon
emissions.  To solve this three-objective problem, a Hybrid Non-dominated
Sorting Genetic Algorithm~III (HNSGA-III) is proposed, incorporating four
enhancements over standard NSGA-III: (1)~nearest-neighbor (NN) initialization
for geographically structured starting populations; (2)~hybrid mutation
combining 3-point inversion and 2-point swap operators; (3)~embedded 2-opt
local search applied to all offspring per generation; and (4)~stagnation-triggered
Variable Neighborhood Search (VNS) using exchange and relocate operators to
escape local optima.  Carbon emissions are quantified using the MEET
(Methodology for calculating Transport Emissions and Energy consumption)
framework~\cite{hickman1999}.  First-echelon routing is solved exactly via a
Gurobi-based MILP solver after the second-echelon routes are optimized by
HNSGA-III.

Experiments on 20 benchmark instances (50-customer and 80-customer groups,
10 independent runs each) compare HNSGA-III against NK-HNSGA-III (a K-Medoids
initialization variant) and standard NSGA-III\@.  HNSGA-III achieves the lowest
average cost (\$10{,}567.5) among all algorithms, outperforming NSGA-III by
1.1\% while reducing solution variance by 51.0\% (Std: 96.5 vs.\ 197.1).
Pareto front quality metrics confirm further advantages: HNSGA-III achieves
superior hypervolume (HV\,=\,0.7995 vs.\ 0.7635), better solution spacing
(SP\,=\,0.1535 vs.\ 0.1924), and a 36.4\% lower generational distance
(GD\,=\,0.1956 vs.\ 0.3076) than NSGA-III\@.  Set coverage analysis shows
that 58.0\% of NSGA-III's Pareto solutions are dominated by HNSGA-III,
compared to 46.0\% in the reverse direction.  Statistically significant cost
improvements (Mann-Whitney U, $p<0.05$) are confirmed in 2 of 20 instances.
These results demonstrate that the proposed HNSGA-III effectively improves
solution quality, consistency, and Pareto front coverage for large-scale green
two-echelon logistics optimization.

\medskip
\noindent\textbf{Keywords:} Two-echelon vehicle routing problem;
Multi-objective optimization; NSGA-III; Nearest-neighbor initialization;
Variable neighborhood search; Carbon emissions; Hypervolume; Green logistics
\end{abstract}

% ============================================================
\section{Introduction}
\label{sec:intro}
% ============================================================

This section establishes the motivation for studying the 2E-MD-MD-MP-CVRP-SPD,
defines the specific problem considered, identifies limitations in existing
solution approaches, and summarizes the contributions and organization of
the paper.

\subsection{Background and Motivation}

The global logistics industry has grown into one of the most critical pillars
of modern economies.  With the global logistics market projected to surpass
USD~12.6 trillion by~2027~\cite{grandview2020}, the pressure to deliver goods
efficiently, cost-effectively, and sustainably has never been greater.  The
rise of e-commerce, the proliferation of urban distribution centers, and
increasingly dynamic customer demand patterns have collectively transformed
last-mile and two-echelon logistics into highly complex optimization challenges.

Vehicle Routing Problems (VRPs) constitute a fundamental class of combinatorial
optimization problems in logistics.  Among their extensions, the two-echelon
vehicle routing problem (2E-VRP) has attracted growing research interest as it
captures the two-tier distribution structure commonly observed in urban freight
systems: a first echelon connects central depots to intermediate facilities
(satellites), while a second echelon handles deliveries from satellites to end
customers.  Incorporating real-world constraints such as multiple depots,
multi-day planning horizons, multiple daily time periods, vehicle capacity
limits, and simultaneous pickup and delivery yields the
2E-MD-MD-MP-CVRP-SPD, a problem formulation that closely reflects operational
realities faced by modern logistics providers.

Beyond economic efficiency, environmental sustainability has become a central
concern for the logistics sector.  Transportation is responsible for
approximately 23\% of global CO$_2$ emissions~\cite{iea2022}, and regulatory
bodies worldwide are tightening emission standards.  Green logistics---the
integration of environmental objectives into route planning---has therefore
emerged as an important research direction.  In this study, carbon emissions
are quantified using the MEET framework~\cite{hickman1999}, and minimizing
carbon output is embedded as a third optimization objective alongside total
cost and vehicle count.

\subsection{Problem Overview}

The 2E-MD-MD-MP-CVRP-SPD considered in this study incorporates the following
real-world features:
\begin{enumerate}
  \item \textit{Two-echelon structure}: First-echelon vehicles transport goods
        from central depots to satellite facilities; second-echelon vehicles
        serve customers from satellites within prescribed time windows.
  \item \textit{Multi-depot}: Multiple central depots and multiple satellite
        locations are available, requiring joint optimization of depot-to-satellite
        and satellite-to-customer assignments.
  \item \textit{Multi-day and multi-period}: The planning horizon spans multiple
        days, and each day is divided into multiple service periods
        (e.g., morning and afternoon), with customer demand varying across both
        days and periods.
  \item \textit{Simultaneous pickup and delivery (SPD)}: Each second-echelon
        vehicle may collect returns from customers while delivering goods on the
        same route, demanding feasible load sequencing.
  \item \textit{Three-objective optimization}: The model simultaneously minimizes
        (i)~total transportation cost, (ii)~total number of vehicles deployed,
        and (iii)~total carbon emissions, forming a true multi-objective problem
        with non-trivial trade-offs.
\end{enumerate}

\subsection{Limitations of Existing Approaches}

Existing solution approaches for 2E-VRP variants predominantly rely on either
exact optimization (e.g., Gurobi-based MILP, branch-and-price) or classical
metaheuristics (e.g., NSGA-II, adaptive large neighborhood search).  While
exact methods guarantee optimality, they scale poorly beyond small instances.
Classical evolutionary algorithms such as NSGA-II and standard NSGA-III often
suffer from:
\begin{enumerate}[label=(\roman*)]
  \item Poor initial solution quality due to random initialization;
  \item Insufficient diversity maintenance in later evolutionary stages;
  \item Weak local refinement capability at the chromosome level;
  \item Limited robustness across independent runs, producing high solution
        variance.
\end{enumerate}

\subsection{Proposed Approach and Contributions}

This study proposes HNSGA-III, which addresses the four limitations above
through four targeted enhancements:
\begin{enumerate}
  \item \textbf{NN Initialization}: Constructs each initial individual as a
        geographically coherent nearest-neighbor tour with a randomly chosen
        starting customer.
  \item \textbf{Hybrid Mutation}: Alternates between 3-point inversion
        (exploration) and 2-point swap (exploitation), each selected with
        equal probability.
  \item \textbf{Embedded 2-opt}: Applied to every offspring before NSGA-III
        selection, raising the average quality of the candidate pool.
  \item \textbf{Stagnation-Triggered VNS}: Exchange and relocate operators
        are applied to the best individual when no improvement occurs for
        $\tau=5$ consecutive generations.
\end{enumerate}

The main contributions are: (1)~a comprehensive three-objective
2E-MD-MD-MP-CVRP-SPD formulation integrating the MEET carbon model;
(2)~the HNSGA-III algorithm with the four enhancements above; (3)~a hybrid
two-tier solution framework coupling HNSGA-III (SE layer) with Gurobi MILP
(FE layer); and (4)~empirical evaluation using six performance metrics across
20 benchmark instances.

\subsection{Paper Organization}

\Cref{sec:lit} reviews related literature.  \Cref{sec:model} presents the
mathematical model.  \Cref{sec:alg} describes HNSGA-III\@.  \Cref{sec:exp}
reports computational experiments.  \Cref{sec:conc} concludes.

% ============================================================
\section{Literature Review}
\label{sec:lit}
% ============================================================

The following subsections survey the principal literature streams relevant to
the 2E-MD-MD-MP-CVRP-SPD: two-echelon routing and its structural extensions,
green VRP with carbon emission modeling, multi-objective evolutionary
algorithms applied to routing, and local search hybridization.  The section
closes by positioning the present work relative to the identified research gap.

\subsection{Two-Echelon VRP and Extensions}

The Vehicle Routing Problem, introduced by Dantzig and
Ramser~\cite{dantzig1959}, is among the most extensively studied problems in
operations research.  The Two-Echelon VRP (2E-VRP), formalized by
Perboli et al.~\cite{perboli2011} and Crainic et al.~\cite{crainic2011},
introduces intermediate satellite facilities between depots and customers.
Sluijk et al.~\cite{sluijk2023} provide a comprehensive survey systematically
categorizing 2E-VRP variants.

The SPD extension requires SE vehicles to collect returns from customers while
delivering goods on the same route.  Belgin et al.~\cite{belgin2018} proposed
one of the first exact models for 2E-CVRP-SPD\@.  Zhou
et al.~\cite{zhou2022spd} extended this with time window constraints, and
Li et al.~\cite{li2022stoch} addressed 2E-VRP-SPD with stochastic demand.

The multi-depot (MD) extension recognizes that real logistics networks feature
multiple warehouses.  Kancharla and Ramadurai~\cite{kancharla2019} formulated
MD-2E-VRP for urban freight, showing that coordinated satellite selection
significantly reduces total cost.  Zhou et al.~\cite{zhou2018md} addressed
MD-2E-VRP with time windows via a hybrid genetic algorithm.

The multi-period dimension extends the planning horizon to multiple days.
Mourgaya and Vanderbeck~\cite{mourgaya2007} proposed a column generation
approach for multi-period VRP\@.  Wen et al.~\cite{wen2010} developed a
multi-period two-echelon model for city logistics.  Wang
et al.~\cite{wang2021} further addressed multi-depot, multi-period 2E-VRP
for cold-chain logistics, and Wang et al.~\cite{wang2023} extended this to
collaborative satellite sharing.  No prior work simultaneously addresses all
five characteristics of the 2E-MD-MD-MP-CVRP-SPD.

\subsection{Green VRP and Carbon Emission Modeling}

Transportation accounts for approximately 23\% of global CO$_2$
emissions~\cite{iea2022}.  The MEET model~\cite{hickman1999} provides a
comprehensive multi-component emission estimation framework, decomposing
vehicle emissions into the base emission rate function $\varphi(v)$, the load
correction function $\psi(v)$, and the road slope correction $r(v)$.
Li et al.~\cite{li2025} applied MEET to a multi-objective green VRP with time
windows, demonstrating that MEET-based objectives produce qualitatively
different Pareto frontiers compared to simple distance-based proxies.

Green 2E-VRP research has grown substantially.  Jepsen
et al.~\cite{jepsen2013} introduced a branch-and-cut exact method for the
symmetric two-echelon capacitated VRP\@.  Multi-objective green VRP
formulations show that Pareto-optimal trade-off frontiers between cost and
emissions are non-trivial and contain solutions qualitatively superior to
single-objective optima.

\subsection{Multi-Objective Evolutionary Algorithms for VRP}

NSGA-II~\cite{deb2002nsga2} introduced fast non-dominated sorting and
crowding-distance-based selection, becoming the de facto reference MOEA for VRP
applications.  NSGA-III~\cite{deb2014nsga3} addresses its limitation in
three-or-more-objective problems by replacing crowding distance with structured
reference points uniformly distributed on a normalized hyperplane.

Standard NSGA-III exhibits weaknesses for VRP: random initialization produces
low-quality starting solutions, standard operators lack route-level insight,
and no explicit local-optima escape mechanism exists.  These weaknesses have
motivated improved NSGA-III variants incorporating transfer learning-based
initialization~\cite{improvedNSGA3transfer}, deep Q-learning guided operator
selection~\cite{dqlNSGA3}, and co-evolutionary strategies~\cite{coevNSGA3}.

Li et al.~\cite{li2025} proposed INSGA-III for the MOGVRPTW\@.  The method
incorporates NN initialization and selective 2-opt on elite individuals,
achieving significant cost reductions over NSGA-II and NSGA-III on Solomon
benchmarks.
The proposed HNSGA-III extends this to the substantially more complex
2E-MD-MD-MP-CVRP-SPD by incorporating hybrid mutation with
stagnation-triggered VNS, applying 2-opt to \emph{all} offspring, and
coupling the SE optimizer with exact FE MILP\@.

\subsection{Local Search and Hybrid Metaheuristics for VRP}

The 2-opt operator~\cite{lin1965} reverses a route segment to eliminate
crossing arcs and remains one of the most computationally efficient route
improvement procedures.  Embedding 2-opt within a genetic algorithm's
generational cycle consistently improves solution quality versus
post-processing-only application~\cite{li2025}.

The Or-opt family~\cite{or1976}, including relocate and exchange operators,
extends 2-opt by allowing non-reversal moves.  The VNS
framework~\cite{hansen2001} alternates between neighborhoods of increasing
size, preventing premature convergence.  Adaptive Large Neighborhood
Search~\cite{ropke2006} provides a general perturbation framework, though its
operator-selection overhead adds implementation complexity relative to the
targeted VNS used in this study.

The k-means clustering heuristic is widely used for customer-satellite
assignment in multi-satellite routing, providing a fast geometric partitioning
consistent with approaches reported for multi-satellite
2E-VRP~\cite{wang2021,wang2023}.

\subsection{Research Gap and Positioning}

No prior work simultaneously addresses: (1)~two-echelon structure with
multiple depots and satellites; (2)~multi-day, multi-period planning with
period-dependent demand; (3)~SPD at the second echelon; (4)~three competing
objectives via MEET; and (5)~a hybrid MOEA combining NN initialization,
hybrid mutation, 2-opt on all offspring, stagnation-triggered VNS, and exact
FE MILP\@.  The 2E-MD-MD-MP-CVRP-SPD and HNSGA-III proposed in this study
close this gap.

% ============================================================
\section{Problem Statement and Mathematical Model}
\label{sec:model}
% ============================================================

This section formally defines the 2E-MD-MD-MP-CVRP-SPD, states the modeling
assumptions, introduces notation, and presents the complete three-objective
mathematical formulation including all first- and second-echelon constraints.

\subsection{Problem Statement}

The 2E-MD-MD-MP-CVRP-SPD is a multi-objective logistics optimization problem
defined over a two-level distribution network.  At the first echelon (FE),
large-capacity trucks transport consolidated freight from depots to satellites.
At the second echelon (SE), smaller vehicles depart from satellites to serve
customers simultaneously with deliveries and pickups.  The planning horizon
spans multiple days ($t\in T$), each divided into multiple service periods
($p\in P$), with customer demand varying across days and periods.

The system operates in three stages: (i)~k-means preprocessing clusters
customers to satellites; (ii)~HNSGA-III determines SE vehicle routes for each
satellite-period block; (iii)~a Gurobi MILP solver optimizes FE routing given
the SE solution.  Three objectives are minimized simultaneously:
$f_1$~(total transportation cost), $f_2$~(total vehicle count), and
$f_3$~(total carbon emissions via MEET).

\subsection{Assumptions}

\begin{enumerate}
  \item All goods collected from customers are temporarily stored at the
        serving satellite before being returned to depots.
  \item Each customer has a soft service time window $[a_i^{t,p},\,b_i^{t,p}]$;
        early or late arrival incurs a penalty cost.
  \item The planning horizon covers multiple operational days, each subdivided
        into multiple periods; customers may appear at different locations on
        different days.
  \item Customer locations, delivery demands, and pickup quantities are known
        for each day and period.
  \item Each depot has sufficient stock capacity to supply all customer demands.
  \item Satellites are not subject to throughput capacity constraints.
  \item The total load (delivery + pickup) on any SE vehicle route must not
        exceed $\QSE$.
  \item Each SE vehicle route must originate from and return to the same
        satellite.
  \item Delivery and pickup operations for each customer are performed
        simultaneously by the same vehicle on the same visit.
  \item SE vehicles cannot serve customers assigned to different satellites
        on a single route.
  \item All FE trucks are homogeneous with maximum load capacity $\QFE$;
        any truck may serve any depot-to-satellite arc.
  \item All SE vehicles are homogeneous with maximum load capacity $\QSE$;
        SE vehicles travel at fixed speed $\vSE = 50$~km/h.
\end{enumerate}

\subsection{Notation}

\paragraph{Sets and indices.}

\begin{tblr}{l l}
  $D,\,d$          & Set of depots; $d$ is the index \\
  $S,\,s$          & Set of satellites; $s$ is the index \\
  $C,\,i,j$        & Set of customers; $i,j$ are indices \\
  $T,\,t$          & Set of planning days; $t$ is the index \\
  $P,\,p$          & Set of service periods; $p$ is the index \\
  $\KFE,\,r$       & Set of FE trucks; $r$ is the index \\
  $\KSE,\,k$       & Set of SE vehicles; $k$ is the index \\
  $\NFE = D\cup S$ & Node set for FE routing \\
  $\NSE = S\cup C$ & Node set for SE routing \\
\end{tblr}

\paragraph{Parameters.}

\begin{tblr}{l l}
  $L_{ij}$              & Distance from node $i$ to node $j$ (km) \\
  $\vSE$                & SE vehicle speed (= 50 km/h) \\
  $\vFE$                & FE vehicle speed (km/h) \\
  $TT_{ds}$             & Travel time from depot $d$ to satellite $s$ \\
  $DD_i^{t,p}$          & Delivery demand of customer $i$ on day $t$, period $p$ (kg) \\
  $DP_i^{t,p}$          & Pickup demand of customer $i$ on day $t$, period $p$ (kg) \\
  $DD_s^t$              & Total delivery demand at satellite $s$ on day $t$ (kg) \\
  $DP_s^t$              & Total pickup demand at satellite $s$ on day $t$ (kg) \\
  $\QSE$                & SE vehicle capacity (= 400 kg) \\
  $\QFE$                & FE vehicle capacity (= 4{,}000 kg) \\
  $\cSE$                & SE unit transportation cost (= 0.05 \$/km) \\
  $\cFE$                & FE unit transportation cost (= 0.2 \$/km) \\
  $F_r$                 & FE truck fixed deployment cost (= 80 \$) \\
  $\xi_e,\,\xi_l$       & Early/late arrival penalty rates (= 0.2 \$/min each) \\
  $a_d,\,b_d$           & Operating time window of depot $d$ \\
  $a_s^{t,p},b_s^{t,p}$ & Time window of satellite $s$ on day $t$, period $p$ \\
  $a_i^{t,p},b_i^{t,p}$ & Service time window of customer $i$ on day $t$, period $p$ \\
\end{tblr}

\paragraph{Decision variables.}

\begin{tblr}{l l}
  $x_{ij}^{k,s,t,p}\in\{0,1\}$           & 1 if SE vehicle $k$ travels arc $(i,j)$ at
                                             satellite $s$, day $t$, period $p$ \\
  $u_i^{k,s,t,p}\ge 0$                   & Cumulative load of SE vehicle $k$ at customer $i$ \\
  $A_i^{k,s,t,p}\ge 0$                   & Arrival time of SE vehicle $k$ at node $i$ \\
  $y_{ds}^{r,t}\in\{0,1\}$               & 1 if FE truck $r$ travels arc $(d,s)$ on day $t$ \\
  $f_{ds}^{r,t}\ge 0$                    & Freight volume on arc $(d,s)$ by truck $r$, day $t$ \\
  $ST_{dr}^t\ge 0$                       & Departure time of FE truck $r$ from depot $d$,
                                            day $t$ \\
  $VT_{sr}^t\ge 0$                       & Arrival time of FE truck $r$ at satellite $s$,
                                            day $t$ \\
  $z_s^{t,p}\in\{0,1\}$                  & 1 if satellite $s$ is open on day $t$, period $p$ \\
  $\delta_r^{\mathrm{FE},t}\in\{0,1\}$   & 1 if FE truck $r$ is deployed on day $t$ \\
  $\delta_k^{\mathrm{SE},t,p}\in\{0,1\}$ & 1 if SE vehicle $k$ is deployed on day $t$,
                                            period $p$ \\
\end{tblr}

\subsection{Mathematical Model}

\subsubsection{Objective Functions}

\paragraph{(f1) Total cost minimization.}
FE transportation cost $\CFE$ includes a distance-proportional term and a
per-truck fixed deployment cost:
\begin{equation}
  \CFE = \sum_{d\in D}\sum_{s\in\NFE}\sum_{r\in\KFE}\sum_{t\in T}
           y_{ds}^{r,t}\,L_{ds}\,\cFE
         + \sum_{r\in\KFE}\sum_{t\in T}\delta_r^{\mathrm{FE},t}\,F_r
  \label{eq:CFE}
\end{equation}
where $\cFE = 0.2$ \$/km and $F_r = 80$ \$.  SE cost $\CSE$ combines
distance-proportional transport and soft time-window penalties:
\begin{equation}
  \CSE = \sum_{i,j\in\NSE}\sum_{k\in\KSE}\sum_{s\in S}\sum_{t\in T}\sum_{p\in P}
           x_{ij}^{k,s,t,p}\,L_{ij}\,\cSE
         + \sum_{i\in C}\sum_{k\in\KSE}\sum_{s\in S}\sum_{t\in T}\sum_{p\in P}
           \Bigl[\xi_e\max\!\bigl(a_i^{t,p}-A_i^{k,s,t,p},0\bigr)
                +\xi_l\max\!\bigl(A_i^{k,s,t,p}-b_i^{t,p},0\bigr)\Bigr]
  \label{eq:CSE}
\end{equation}
where $\cSE=0.05$ \$/km; $\xi_e=\xi_l=0.2$ \$/min; morning vehicles depart
satellite at $t=540$ min (09:00), afternoon at $t=720$ min (12:00).
\begin{equation}
  \min\;f_1 = \CFE + \CSE
  \label{eq:f1}
\end{equation}

\paragraph{(f2) Total vehicle count minimization.}
\begin{equation}
  \min\;f_2 = \sum_{t\in T}\sum_{p\in P}
    \Bigl(\sum_{r\in\KFE}\delta_r^{\mathrm{FE},t}
         +\sum_{k\in\KSE}\delta_k^{\mathrm{SE},t,p}\Bigr)
  \label{eq:f2}
\end{equation}

\paragraph{(f3) Total carbon emissions minimization.}
Carbon emissions are computed using the MEET model~\cite{hickman1999}.
For a 3.5-ton box truck (Table~A24 of~\cite{hickman1999}):
\begin{align}
  \varphi(v) &= 110 + 0.000375\,v^3 + \frac{8702}{v}
  \quad\text{(g CO}_2\text{/km)}
  \label{eq:phi} \\
  \psi(v)    &= 1.27 - 0.00235\,v - \frac{1.33}{v}
  \label{eq:psi}
\end{align}
The unit carbon emission is $\varphi(v)\cdot\psi(v)/1000$ (kg CO$_2$/km).
At the fixed SE speed $\vSE=50$ km/h:
\[
  \varphi(50)=330.915\;\text{g/km},\quad
  \psi(50)=1.1259,\quad
  \eSE=330.915\times1.1259/1000\approx0.3726\;\text{kg CO}_2\text{/km.}
\]
\begin{align}
  \ESE &= \sum_{i,j\in\NSE}\sum_{k\in\KSE}\sum_{s\in S}\sum_{t\in T}\sum_{p\in P}
           x_{ij}^{k,s,t,p}\,L_{ij}\,\eSE
  \label{eq:ESE}\\
  \EFE &= \sum_{(d,s)}\sum_{r\in\KFE}\sum_{t\in T}
           y_{ds}^{r,t}\,L_{ds}\,e_{\mathrm{FE}}(\vFE)
  \label{eq:EFE}
\end{align}
\begin{equation}
  \min\;f_3 = \ESE + \EFE
  \label{eq:f3}
\end{equation}

\subsubsection{First-Echelon Constraints}

\begin{align}
  &\sum_{d\in D}\sum_{r\in\KFE}y_{ds}^{r,t} \ge z_s^{t,p},
  &&\forall\,s\in S,\,t\in T,\,p\in P
  \tag{C1}\label{c1}\\
  &\sum_{d\in D,s\in S}y_{ds}^{r,t}-\sum_{s\in S,d\in D}y_{sd}^{r,t}=0,
  &&\forall\,r\in\KFE,\,t\in T
  \tag{C2}\label{c2}\\
  &\sum_{d\in D,s\in S}y_{ds}^{r,t}(DD_s^t+DP_s^t)\le\QFE,
  &&\forall\,r\in\KFE,\,t\in T
  \tag{C3}\label{c3}\\
  &f_{ds}^{r,t}\ge 0,
  &&\forall\,d\in D,s\in S,r\in\KFE,t\in T
  \tag{C4a}\label{c4a}\\
  &\sum_{d\in D}f_{ds}^{r,t}-\sum_{d\in D}f_{sd}^{r,t}
    =y_{ds}^{r,t}(DD_s^t+DP_s^t),
  &&\forall\,s\in S,r\in\KFE,t\in T
  \tag{C4b}\label{c4b}\\
  &z_s^{t,p}\le\sum_{d\in D}\sum_{r\in\KFE}y_{ds}^{r,t},
  &&\forall\,s\in S,t\in T,p\in P
  \tag{C5}\label{c5}\\
  &a_d\sum_{s\in S}y_{ds}^{r,t}\le ST_{dr}^t\le b_d\sum_{s\in S}y_{ds}^{r,t},
  &&\forall\,d\in D,r\in\KFE,t\in T
  \tag{C6}\label{c6}\\
  &ST_{dr}^t+TT_{ds}-M(1-y_{ds}^{r,t})\le VT_{sr}^t
    \le ST_{dr}^t+TT_{ds}+M(1-y_{ds}^{r,t}),
  &&\forall\,d,s,r\in\KFE,t
  \tag{C7}\label{c7}\\
  &y_{ds}^{r,t}\in\{0,1\},\quad
  f_{ds}^{r,t}\ge 0,\quad
  z_s^{t,p}\in\{0,1\},
  &&\forall\,d,s,r,t,p
  \tag{C8}\label{c8}
\end{align}

\subsubsection{Second-Echelon Constraints}

\begin{align}
  &\sum_{j\in\NSE}\sum_{k\in\KSE}x_{ij}^{k,s,t,p}=1,
  &&\forall\,i\in C,s\in S,t\in T,p\in P
  \tag{C9}\label{c9}\\
  &\sum_{j\in\NSE}x_{ij}^{k,s,t,p}-\sum_{j\in\NSE}x_{ji}^{k,s,t,p}=0,
  &&\forall\,i\in\NSE,k\in\KSE,s,t,p
  \tag{C10}\label{c10}\\
  &\sum_{i\in C,j\in\NSE}x_{ij}^{k,s,t,p}(DD_i^{t,p}+DP_i^{t,p})\le\QSE,
  &&\forall\,k\in\KSE,s\in S,t\in T,p\in P
  \tag{C11}\label{c11}\\
  &\sum_{j\in C}x_{sj}^{k,s,t,p}=\sum_{j\in C}x_{js}^{k,s,t,p}\le 1,
  &&\forall\,k\in\KSE,s\in S,t\in T,p\in P
  \tag{C12}\label{c12}\\
  &u_i^{k,s,t,p}-u_j^{k,s,t,p}+\QSE\,x_{ij}^{k,s,t,p}\le\QSE-DD_j^{t,p},
  &&\forall\,i,j\in C,k,s,t,p
  \tag{C13a}\label{c13a}\\
  &DD_i^{t,p}\le u_i^{k,s,t,p}\le\QSE,
  &&\forall\,i\in C,k,s,t,p
  \tag{C13b}\label{c13b}\\
  &A_j^{k,s,t,p}\ge A_i^{k,s,t,p}+\frac{L_{ij}}{\vSE}
    -M(1-x_{ij}^{k,s,t,p}),
  &&\forall\,i,j\in C,k,s,t,p
  \tag{C14}\label{c14}\\
  &a_s^{t,p}\sum_{j\in C}x_{sj}^{k,s,t,p}\le A_s^{k,s,t,p}
    \le b_s^{t,p}\sum_{j\in C}x_{sj}^{k,s,t,p},
  &&\forall\,k\in\KSE,s,t,p
  \tag{C15}\label{c15}\\
  &x_{ij}^{k,s,t,p}\in\{0,1\},\quad
  u_i^{k,s,t,p}\ge 0,\quad
  A_i^{k,s,t,p}\ge 0.
  \tag{C16}\label{c16}
\end{align}

\subsection{Model Summary}

The complete formulation is $\min\;\{f_1,f_2,f_3\}$ subject to
\eqref{c1}--\eqref{c16}.  The problem is NP-hard (even single-echelon CVRP
is NP-hard); the HNSGA-III + Gurobi decomposition described in \cref{sec:alg}
provides a practically scalable solution approach.

% ============================================================
\section{Proposed Algorithm: HNSGA-III}
\label{sec:alg}
% ============================================================

This section describes HNSGA-III in detail, covering solution encoding,
the four enhancements over standard NSGA-III, and the complete algorithmic
procedure including the Gurobi-based FE solver.

\subsection{Overview}

HNSGA-III builds upon the standard NSGA-III reference-point
framework~\cite{deb2014nsga3} and introduces four targeted enhancements
described in Sections~\ref{sec:init}--\ref{sec:vns}.  The algorithm operates
on the SE routing layer; FE routing is solved exactly by Gurobi MILP after SE
routes are finalized.

\subsection{Solution Encoding and Parameters}

Each chromosome is a permutation of customer IDs assigned to a given satellite
in a given time period.  For $n_s$ satellites and two daily periods,
$2n_s$ independent populations of size $N_{\mathrm{pop}}=30$ are maintained.
\Cref{tab:params} summarizes all algorithm parameters.

\begin{table}[ht]
\centering
\caption{Algorithm Parameter Settings}
\label{tab:params}
\begin{tblr}{
    colspec = {l c},
}
\hline[0.08em]
Parameter & Value \\ \hline[0.05em]
Population size ($N_{\mathrm{pop}}$)              & 30 \\
Number of generations                             & 1{,}000 \\
Crossover probability ($cp$) --- Instances 1--10  & 0.70 \\
Crossover probability ($cp$) --- Instances 11--20 & 0.55 \\
Mutation probability ($mp$) --- Instances 1--10   & 0.70 \\
Mutation probability ($mp$) --- Instances 11--20  & 0.80 \\
SE vehicle capacity $\QSE$                        & 400 kg \\
SE vehicle speed $\vSE$                           & 50 km/h \\
FE vehicle capacity $\QFE$                        & 4{,}000 kg \\
Stagnation threshold $\tau$ (VNS trigger)         & 5 generations \\
Independent runs per instance                     & 10 \\ \hline[0.08em]
\end{tblr}
\end{table}

\subsection{Population Initialization: Nearest-Neighbor Heuristic}
\label{sec:init}

For each individual, a random customer is selected as the starting point.
The algorithm iteratively appends the nearest unvisited customer (by Euclidean
distance) until all customers are visited, forming an NN tour.  Because each
individual begins from a different starting customer, the $N_{\mathrm{pop}}$
individuals collectively cover diverse geographic entry points while
individually maintaining local proximity structure.  \Cref{alg:nn} describes the procedure.

\begin{algorithm}[ht]
\caption{NN Initialization for one population}
\label{alg:nn}
\begin{algorithmic}[1]
\Require Customer set $C$, coordinate lookup \textit{coord}, target size
         $N_{\mathrm{pop}}$
\Ensure Initial population $P$
\State $P \leftarrow \emptyset$;\quad \textit{attempts} $\leftarrow 0$
\While{$|P| < N_{\mathrm{pop}}$ \textbf{and} \textit{attempts}
       $< 30\cdot N_{\mathrm{pop}}$}
  \State $\textit{start} \leftarrow \mathrm{random}(C)$
  \State $\textit{seq} \leftarrow [\textit{start}]$;\quad
         $\textit{unvisited} \leftarrow C \setminus \{\textit{start}\}$
  \While{$\textit{unvisited} \neq \emptyset$}
    \State $\textit{next} \leftarrow
           \arg\min_{u\in\textit{unvisited}}\mathrm{dist}(\textit{seq.last},u)$
    \State $\textit{seq}.\mathrm{append}(\textit{next})$;\quad
           $\textit{unvisited} \leftarrow \textit{unvisited}\setminus
           \{\textit{next}\}$
  \EndWhile
  \If{$\textit{seq} \notin P$}
    \State $P \leftarrow P \cup \{\textit{seq}\}$
  \EndIf
  \State $\textit{attempts} \leftarrow \textit{attempts}+1$
\EndWhile
\While{$|P| < N_{\mathrm{pop}}$}
  \State $P \leftarrow P \cup \{\mathrm{random\_shuffle}(C)\}$
\EndWhile
\State \Return $P$
\end{algorithmic}
\end{algorithm}

\subsection{Crossover: Order-Based Crossover (OBX)}

OBX selects a random contiguous segment $[a,b]$ from parent $p_1$, copies it
to offspring $c_1$ at the same positions, then fills remaining positions with
customers from $p_2$ in the order they appear, skipping already-placed
customers.  Offspring $c_2$ is generated symmetrically.

\subsection{Hybrid Mutation Operator}

Two mutation operators are defined, each selected with equal probability (50\%)
when mutation is triggered:
\begin{itemize}
  \item \textbf{3-Point Inversion (exploration)}: Three positions $i<j<k$ are
        selected; the subsequence $[i,k]$ is reversed.  Produces macro-level
        structural changes.
  \item \textbf{2-Point Swap (exploitation)}: Two positions $a$ and $b$ are
        selected and their customers are exchanged.  Produces fine-grained
        local adjustment.
\end{itemize}

\subsection{Selection: NSGA-III Reference-Point Framework}

The combined parent--offspring pool is reduced to $N_{\mathrm{pop}}$ using
standard NSGA-III~\cite{deb2014nsga3}: (1)~non-dominated sorting partitions
individuals into fronts $F_1,F_2,\ldots$; (2)~fronts are admitted until the
population would exceed $N_{\mathrm{pop}}$; (3)~reference points on the
normalized hyperplane govern niching selection in the critical last front.

\subsection{Embedded 2-opt Local Search}

After crossover and mutation, a 2-opt improvement is applied to every
offspring: decode the chromosome into routes, select one route at random,
reverse a subsequence $[a,b]$ if it yields a strictly better non-duplicate
solution.  Applying 2-opt to \emph{all} offspring (rather than only elites)
ensures each candidate receives route-level refinement before competing in
selection.

\subsection{Stagnation-Triggered Variable Neighborhood Search}
\label{sec:vns}

A stagnation counter per satellite-period block is incremented every
generation in which the best-individual cost fails to improve.  When the counter reaches
$\tau=5$, a two-step VNS is applied to the current best individual:
\begin{enumerate}
  \item \textbf{Exchange}: two customers within the same route are swapped.
  \item \textbf{Relocate}: one customer is removed and reinserted at a
        different position.
\end{enumerate}
If either step produces a strictly better non-duplicate solution, it replaces
the current best and the stagnation counter resets.

\subsection{Complete Algorithm}

\Cref{alg:hnsga3} presents the complete HNSGA-III procedure.

\begin{algorithm}[ht]
\caption{HNSGA-III for 2E-MD-MD-MP-CVRP-SPD}
\label{alg:hnsga3}
\begin{algorithmic}[1]
\Require Instance data, parameters ($N_{\mathrm{pop}},cp,mp,\tau,
         \mathrm{num\_gen}$)
\Ensure Best SE routes per satellite-period; FE cost
\State Cluster customers to satellites via K-Means
\State Partition customers into morning / afternoon periods
\For{each satellite $s$ and period $p$}
  \State $P_{s,p} \leftarrow \mathrm{NN\_Initialize}(\mathrm{customers}_{s,p},
         N_{\mathrm{pop}})$
         \quad\Comment{see \cref{alg:nn}}
\EndFor
\State $\mathrm{stag}_{s,p} \leftarrow 0$ for all $s,p$
\For{$\mathrm{gen} = 1$ \To $\mathrm{num\_gen}$}
  \For{each population $P_{s,p}$}
    \State $\mathrm{best\_before} \leftarrow \min\mathrm{cost}(P_{s,p})$
    \If{$\mathrm{rand}() < cp$}
      \State $O \leftarrow \mathrm{OBX}(P_{s,p})$
    \EndIf
    \State $O \leftarrow O \cup \mathrm{HybridMutate}(P_{s,p},mp)$
    \State $O \leftarrow \mathrm{2\text{-}opt}(O)$
    \State $P_{s,p} \leftarrow \mathrm{NSGA3\_Select}(P_{s,p}\cup O,
           N_{\mathrm{pop}})$
    \State $\mathrm{best\_after} \leftarrow \min\mathrm{cost}(P_{s,p})$
    \If{$\mathrm{best\_after} < \mathrm{best\_before}$}
      \State $\mathrm{stag}_{s,p} \leftarrow 0$
    \Else
      \State $\mathrm{stag}_{s,p} \leftarrow \mathrm{stag}_{s,p}+1$
    \EndIf
    \If{$\mathrm{stag}_{s,p} \ge \tau$}
      \State $P_{s,p} \leftarrow \mathrm{VNS}(P_{s,p})$;\quad
             $\mathrm{stag}_{s,p} \leftarrow 0$
    \EndIf
  \EndFor
\EndFor
\State Extract best solution from each $P_{s,p}$
\State Compute satellite demands from SE routes
\State $\mathrm{FE\_cost} \leftarrow
       \mathrm{Gurobi\_MILP}(\mathrm{depots},\mathrm{satellites},
       \mathrm{demands})$
\State \Return SE routes, FE cost
\end{algorithmic}
\end{algorithm}

% ============================================================
\section{Computational Experiments}
\label{sec:exp}
% ============================================================

This section describes the experimental setup, benchmark instances, compared
algorithms, and parameter settings, followed by four complementary performance
analyses covering sensitivity, cost comparison, Pareto front quality, and
statistical significance.

\subsection{Experimental Setup}

\subsubsection{Benchmark Instances}

Twenty benchmark instances are used: Instances~1--10 have 50 customers,
2 satellites, and 1 depot; Instances~11--20 have 80 customers, 3 satellites,
and 1 depot.  Customer coordinates, demands, and soft service-time windows
are generated synthetically over a multi-day, two-period horizon.

\subsubsection{Compared Algorithms}

Three algorithms are compared:
\begin{enumerate}
  \item \textbf{HNSGA-III} (Proposed): full algorithm as described in
        \cref{sec:alg}.
  \item \textbf{NK-HNSGA-III} (Ablation Baseline): HNSGA-III with NN
        initialization replaced by K-Medoids clustering; all other components
        are identical.
  \item \textbf{NSGA-III} (Standard Baseline): reference-point NSGA-III
        \cite{deb2014nsga3} with random initialization and no local search.
\end{enumerate}

\subsubsection{Implementation and Parameter Settings}

All algorithms are implemented in Python~3.  FE routing is solved by Gurobi
MILP solver (version~10.0).  All experiments are conducted on an Intel Core i7-11800H @
2.30~GHz (8 cores / 16 threads) with 32~GB RAM and an NVIDIA GeForce RTX~3070
Laptop GPU\@.  Each algorithm is executed 10 independent times per instance.
Parameters follow \cref{tab:params}; NSGA-III uses $N_{\mathrm{pop}}=30$ and
1{,}000 generations.

\subsection{Parameter Sensitivity Analysis}

\Cref{tab:sensitivity} reports the average best cost over 10 runs for two
representative instances while sweeping crossover probability~$cp$ and
mutation probability~$mp$.

\begin{table}[ht]
\centering
\caption{Parameter Sensitivity Analysis of HNSGA-III}
\label{tab:sensitivity}
\small
\begin{tblr}{
    colspec = {c r r r r r r r r},
}
\hline[0.08em]
\SetCell[r=2]{c} Scenario
  & \SetCell[c=4]{c} Test~A (50 customers, Inst.~1) & & &
  & \SetCell[c=4]{c} Test~B (80 customers, Inst.~11) & & & \\
\cline{2-5}\cline{6-9}
  & $cp$ & Cost(\$) & $mp$ & Cost(\$)
  & $cp$ & Cost(\$) & $mp$ & Cost(\$) \\ \hline[0.05em]
1  & 0.85 & 6339.1 & 0.1 & 6340.1 & 0.85 & 9315.9 & 0.1 & 9406.9 \\
2  & 0.80 & 6327.1 & 0.2 & 6341.1 & 0.80 & 9290.9 & 0.2 & 9421.9 \\
3  & 0.75 & 6339.1 & 0.3 & 6326.1 & 0.75 & 9310.9 & 0.3 & 9315.9 \\
4  & \textbf{*0.70} & \textbf{*6325.1} & 0.4 & 6326.1
   & 0.70 & 9328.9 & 0.4 & 9267.9 \\
5  & 0.65 & 6326.1 & 0.5 & 6337.1 & 0.65 & 9322.9 & 0.5 & 9352.9 \\
6  & 0.60 & 6328.1 & 0.6 & 6327.1 & 0.60 & 9320.9 & 0.6 & 9355.9 \\
7  & 0.55 & 6325.1 & \textbf{*0.7} & \textbf{*6325.1}
   & \textbf{*0.55} & \textbf{*9253.9} & 0.7 & 9354.9 \\
8  & 0.50 & 6326.1 & 0.8 & 6339.1 & 0.50 & 9267.9
   & \textbf{*0.8} & \textbf{*9253.9} \\
9  & 0.45 & 6337.1 & 0.9 & 6325.1 & 0.45 & 9377.9 & 0.9 & 9322.9 \\
10 & 0.40 & 6326.1 & 1.0 & 6337.1 & 0.40 & 9326.9 & 1.0 & 9369.9 \\
\hline[0.08em]
\SetCell[c=9]{l} \small * denotes best setting (lowest cost) in each column. \\
\end{tblr}
\end{table}

\noindent The selected parameters are: Instances~1--10: $cp=0.70$, $mp=0.70$;
Instances~11--20: $cp=0.55$, $mp=0.80$.

\subsection{Algorithm Performance Comparison}

\subsubsection{HNSGA-III Detailed Results}

\Cref{tab:detail} presents per-instance results.  CO$_2$\,avg reports the
total CO$_2$ accumulated across all SE/FE routes (kg), computed via MEET by
multiplying $\varphi(v)\cdot\psi(v)$ (g/km) by route distance and dividing
by 1000 to convert to kilograms.

\begin{table}[ht]
\centering
\caption{HNSGA-III Detailed Results (10 runs per instance)}
\label{tab:detail}
\small
\begin{tblr}{
    colspec = {c r r r r r r r},
}
\hline[0.08em]
Inst & $cp$ & $mp$ & Cost$_{\min}$(\$) & Cost$_{\text{avg}}$(\$)
     & NV$_{\text{avg}}$ & CO$_2$\,avg (kg) & RT$_{\text{avg}}$ (s) \\
\hline[0.05em]
 1 & 0.70 & 0.70 &  6327.1 &  6349.8 & 6.0 & 2662.99 &  53.8 \\
 2 & 0.70 & 0.70 &  6243.3 &  6333.7 & 5.0 & 2550.92 &  63.8 \\
 3 & 0.70 & 0.70 &  8013.0 &  8137.8 & 5.0 & 4417.07 &  61.0 \\
 4 & 0.70 & 0.70 &  6629.9 &  6643.0 & 5.0 & 2286.95 &  54.2 \\
 5 & 0.70 & 0.70 &  5704.5 &  5787.4 & 6.0 & 2272.82 &  55.6 \\
 6 & 0.70 & 0.70 & 11517.1 & 11640.7 & 5.0 & 5206.60 &  62.9 \\
 7 & 0.70 & 0.70 & 11438.5 & 11521.9 & 5.0 & 5896.72 &  60.0 \\
 8 & 0.70 & 0.70 &  7116.5 &  7286.0 & 6.0 & 3467.53 &  52.4 \\
 9 & 0.70 & 0.70 &  8994.0 &  9196.6 & 5.0 & 4290.86 &  61.7 \\
10 & 0.70 & 0.70 &  9266.1 &  9382.8 & 5.0 & 3874.07 &  57.6 \\
\hline[0.05em]
11 & 0.55 & 0.80 &  9342.9 &  9425.1 & 7.0 & 3470.77 &  89.7 \\
12 & 0.55 & 0.80 & 12830.9 & 12873.9 & 8.0 & 4947.96 &  90.7 \\
13 & 0.55 & 0.80 & 12995.9 & 13239.4 & 8.0 & 5001.63 &  89.3 \\
14 & 0.55 & 0.80 & 14124.6 & 14226.6 & 8.0 & 6413.53 & 101.3 \\
15 & 0.55 & 0.80 & 12322.7 & 12526.1 & 8.0 & 4827.86 &  89.7 \\
16 & 0.55 & 0.80 & 15255.4 & 15332.2 & 7.0 & 7139.80 & 102.1 \\
17 & 0.55 & 0.80 & 13335.8 & 13619.2 & 8.0 & 5308.91 &  82.8 \\
18 & 0.55 & 0.80 & 14211.5 & 14308.6 & 8.0 & 5403.76 &  93.2 \\
19 & 0.55 & 0.80 & 13839.3 & 13976.2 & 8.0 & 5403.88 &  98.4 \\
20 & 0.55 & 0.80 &  9863.0 & 10103.5 & 7.0 & 3861.22 &  86.4 \\
\hline[0.05em]
\textbf{Avg} & & & & \textbf{10595.5} & \textbf{6.5}
             & \textbf{4359.38} & \textbf{75.3} \\ \hline[0.08em]
\end{tblr}
\end{table}

\subsubsection{Three-Algorithm Cost Comparison}

\Cref{tab:comparison} compares average cost and standard deviation per
instance across all three algorithms.

\begin{table}[ht]
\centering
\caption{Cost Comparison: HNSGA-III vs NK-HNSGA-III vs NSGA-III}
\label{tab:comparison}
\small
\begin{tblr}{
    colspec = {c r r r r r r r r r},
    colsep = 3pt,
}
\hline[0.08em]
\SetCell[r=2]{c} Inst
  & \SetCell[c=3]{c} HNSGA-III & &
  & \SetCell[c=3]{c} NK-HNSGA-III & &
  & \SetCell[c=3]{c} NSGA-III & & \\
\cline{2-4}\cline{5-7}\cline{8-10}
  & Best & Avg & Std & Best & Avg & Std & Best & Avg & Std \\
\hline[0.05em]
 1 &  6339.1 &  6340.9 &   2.5 &  6337.1 &  6395.7 &  65.3 &  6339.1 &  6416.1 &  73.7 \\
 2 &  6244.3 &  6297.7 &  67.2 &  6217.3 &  6300.9 &  91.3 &  6343.3 &  6382.3 &  27.5 \\
 3 &  8028.0 &  8123.6 &  68.5 &  8016.0 &  8134.6 & 231.2 &  8007.0 &  8051.4 &  58.9 \\
 4 &  6629.9 &  6639.1 &   8.2 &  6371.9 &  6587.1 & 120.5 &  6289.9 &  6439.9 & 193.8 \\
 5 &  5666.5 &  5752.5 &  56.4 &  5717.5 &  5804.5 &  49.0 &  5731.5 &  5812.9 & 113.8 \\
 6 & 11596.1 & 11677.3 & 134.7 & 11593.1 & 11611.7 &  23.2 & 11593.1 & 11685.9 & 181.8 \\
 7 & 11299.5 & 11429.3 &  90.4 & 11299.5 & 11390.1 &  51.2 & 11534.5 & 11617.5 &  49.2 \\
 8 &  7109.5 &  7225.5 & 136.2 &  7111.5 &  7275.7 & 148.9 &  7099.5 &  7362.3 & 249.9 \\
 9 &  8994.0 &  9168.8 & 220.9 &  8996.0 &  9212.0 & 205.7 &  9010.0 &  9109.0 & 139.6 \\
10 &  9285.1 &  9390.3 & 135.0 &  9266.1 &  9360.9 & 113.1 &  9285.1 &  9339.9 &  54.9 \\
\hline[0.05em]
11 &  9244.9 &  9360.3 &  77.8 &  9218.9 &  9376.7 &  90.5 &  9268.9 &  9507.7 & 238.2 \\
12 & 12830.9 & 12900.5 & 124.4 & 12876.9 & 12920.5 &  76.7 & 12853.9 & 13060.7 & 242.4 \\
13 & 13039.9 & 13265.7 & 141.7 & 12660.9 & 13047.5 & 330.6 & 12765.9 & 13324.5 & 508.8 \\
14 & 14124.6 & 14152.6 &  36.2 & 14046.6 & 14212.0 & 124.1 & 14017.6 & 14488.6 & 347.6 \\
15 & 12405.7 & 12502.3 &  68.6 & 12333.7 & 12518.9 & 168.9 & 12488.7 & 12732.9 & 163.2 \\
16 & 15261.4 & 15323.2 &  45.5 & 15212.4 & 15255.6 &  41.8 & 15286.4 & 15586.2 & 331.0 \\
17 & 13396.8 & 13556.4 & 234.9 & 13390.8 & 13689.6 & 354.4 & 13390.8 & 13852.2 & 288.1 \\
18 & 14063.5 & 14174.7 & 107.7 & 13987.5 & 14122.1 & 183.7 & 14014.5 & 14196.7 & 224.5 \\
19 & 13829.3 & 13953.3 &  78.4 & 13940.3 & 14000.9 &  48.2 & 14183.3 & 14437.1 & 307.4 \\
20 &  9990.0 & 10116.0 &  95.5 & 10089.0 & 10160.0 &  74.5 & 10029.0 & 10245.0 & 147.2 \\
\hline[0.05em]
\textbf{Avg} & & \textbf{10567.5} & \textbf{96.5}
             & & \textbf{10568.8} & \textbf{129.6}
             & & \textbf{10682.4} & \textbf{197.1} \\ \hline[0.08em]
\end{tblr}
\end{table}

HNSGA-III achieves the lowest average cost (10{,}567.5), outperforming
NSGA-III by 1.1\%.  The standard deviation is reduced by 51.0\%
(96.5 vs.\ 197.1), and HNSGA-III outperforms NSGA-III on average cost
in 16 of 20 instances.

\subsection{Pareto Front Quality Metrics}

\Cref{tab:pareto} compares six quality metrics averaged across all
20 instances.

\begin{table}[ht]
\centering
\caption{Pareto Front Quality Metrics (averages across 20 instances)}
\label{tab:pareto}
\begin{tblr}{
    colspec = {l r r r r r r},
}
\hline[0.08em]
Algorithm & Avg Cost & Avg Std & Avg SP$\downarrow$
          & Avg HV$\uparrow$ & Avg GD$\downarrow$ & Avg IGD$\downarrow$ \\
\hline[0.05em]
HNSGA-III    & 10{,}567.5 &  96.5 & 0.1535 & 0.7995 & 0.1956 & 0.2174 \\
NK-HNSGA-III & 10{,}568.8 & 129.6 & 0.1549 & 0.7540 & 0.2211 & 0.2018 \\
NSGA-III     & 10{,}682.4 & 197.1 & 0.1924 & 0.7635 & 0.3076 & 0.2087 \\ \hline[0.08em]
\end{tblr}
\end{table}

HNSGA-III achieves superior performance on five of six metrics relative to
NSGA-III\@.  The HV improvement (+4.7\%) confirms broader Pareto front
coverage; the GD improvement ($-36.4\%$) shows solutions are consistently
closer to the best-known frontier.  Comparing HNSGA-III with NK-HNSGA-III
isolates the effect of NN initialization: despite nearly identical average
cost (difference $<0.01\%$), HNSGA-III achieves substantially better
stability (Std: 96.5 vs.\ 129.6), hypervolume (0.7995 vs.\ 0.7540), and
generational distance (0.1956 vs.\ 0.2211).

\subsection{Statistical Analysis}

Mann-Whitney U tests (two-sided, $\alpha=0.05$) are applied to the 10 cost
samples of HNSGA-III and NSGA-III per instance.

\begin{table}[ht]
\centering
\caption{Mann-Whitney U Test: HNSGA-III vs NSGA-III (significant instances only)}
\label{tab:stat}
\begin{tblr}{
    colspec = {c r l l},
}
\hline[0.08em]
Instance & $p$-value & Significant? & Direction \\ \hline[0.05em]
7          & 0.0119  & *            & HNSGA-III better \\
19         & 0.0079  & *            & HNSGA-III better \\
All others & $>0.05$ & ---          & --- \\ \hline[0.08em]
\SetCell[c=4]{l} \small * $p<0.05$ \\
\end{tblr}
\end{table}

Statistically significant improvements are confirmed in 2 of 20 instances.
The Set Coverage (SC) metric shows that 58.0\% of NSGA-III's non-dominated
solutions are covered by HNSGA-III, versus only 46.0\% in the reverse
direction, a directional advantage of 12 percentage points.

\subsection{Satellite Configuration Analysis}

\Cref{tab:satellite} reports the minimum-cost satellite count (1, 2, or
3 satellites open) for each of 30 planning days.

\begin{table}[ht]
\centering
\caption{Optimal Number of Active Satellites per Planning Day (HNSGA-III)}
\label{tab:satellite}
\small
\begin{tblr}{
    colspec = {c r r r r},
}
\hline[0.08em]
Day & Sat=1 (\$) & Sat=2 (\$) & Sat=3 (\$) & Optimal \\ \hline[0.05em]
 1 &  *4324.7 &  5102.0 &  4559.8 & 1 \\
 2 &  7352.5  & *5757.8 &  6765.9 & 2 \\
 3 &  4896.0  &  4900.1 & *4304.4 & 3 \\
 4 &  5180.2  & *4805.3 &  5759.6 & 2 \\
 5 & *5121.7  &  6030.5 &  5808.1 & 1 \\
 6 &  6076.7  & *5478.1 &  6416.4 & 2 \\
 7 &  5124.2  &  6289.5 & *4825.9 & 3 \\
 8 &  7216.9  &  8612.8 & *6925.1 & 3 \\
 9 & *7005.7  &  8339.8 &  8183.2 & 1 \\
10 &  6351.5  & *5512.8 &  6569.3 & 2 \\
11 & *9099.7  & 10495.9 &  9859.7 & 1 \\
12 &  8751.5  & 10902.7 & *8388.4 & 3 \\
13 &  8988.2  & *8057.3 &  9581.8 & 2 \\
14 &  6035.9  & *5970.8 &  7857.3 & 2 \\
15 &  7445.9  &  8983.5 & *6771.6 & 3 \\
16 & *8945.7  & 10201.6 &  9103.1 & 1 \\
17 & *9379.0  & 11229.5 &  9532.3 & 1 \\
18 &  6913.1  &  6784.3 & *6590.4 & 3 \\
19 &  8935.2  & 10130.7 & *8599.5 & 3 \\
20 &*11379.6  & 14687.6 & 12519.2 & 1 \\
21 &*11080.8  & 13100.5 & 14045.3 & 1 \\
22 & 11021.5  & 11733.9 & *8685.9 & 3 \\
23 &*11269.2  & 12508.1 & 12215.5 & 1 \\
24 & 13387.6  & 14106.5 & *9477.4 & 3 \\
25 &  8007.1  &  7175.1 & *7097.0 & 3 \\
26 &*10667.8  & 11552.2 & 11817.6 & 1 \\
27 &*12104.1  & 13799.2 & 13255.6 & 1 \\
28 & 11929.6  & 12271.2 & *8499.7 & 3 \\
29 & *9924.8  & 12867.8 & 10755.7 & 1 \\
30 &*11938.9  & 13215.1 & 12532.9 & 1 \\
\hline[0.05em]
\SetCell[c=5]{l} \textbf{Distribution:} Sat=1: 13 days (43\%),
  Sat=2: 6 days (20\%), Sat=3: 11 days (37\%) \\
\hline[0.08em]
\SetCell[c=5]{l} \small * denotes minimum cost for that day. \\
\end{tblr}
\end{table}

The optimal number of active satellites varies considerably: Sat=1 is optimal
on 43\% of days, Sat=3 on 37\%, and Sat=2 on 20\%, demonstrating that the
optimal satellite configuration is strongly demand-dependent and cannot be
fixed a priori.

% ============================================================
\section{Conclusion}
\label{sec:conc}
% ============================================================

This section summarizes the contributions of the study, discusses the
limitations of the proposed framework, and outlines promising directions for
future research.

\subsection{Summary of Contributions}

This study addressed the 2E-MD-MD-MP-CVRP-SPD, simultaneously minimizing
total transportation cost ($f_1$), total vehicle count ($f_2$), and total
carbon emissions ($f_3$).  HNSGA-III was proposed with four targeted
enhancements over standard NSGA-III: NN initialization, hybrid mutation
(3-point inversion + 2-point swap), embedded 2-opt on all offspring, and
stagnation-triggered VNS ($\tau=5$).  Carbon emissions were quantified
using the MEET model~\cite{hickman1999}; FE routing was solved exactly by
Gurobi MILP after SE routing was optimized by HNSGA-III\@.

Key experimental findings from 20 benchmark instances:
\begin{itemize}
  \item \textbf{Cost and stability}: HNSGA-III achieves the lowest average
        cost (\$10{,}567.5), outperforming NSGA-III by 1.1\% and reducing
        solution variance by 51.0\% (Std: 96.5 vs.\ 197.1).
  \item \textbf{Pareto front quality}: Superior HV ($+4.7\%$), better SP
        ($-20.2\%$), and 36.4\% lower GD than NSGA-III, confirming broader
        and better-distributed frontier coverage.
  \item \textbf{Dominance coverage}: 58.0\% of NSGA-III's non-dominated
        solutions are covered by HNSGA-III vs.\ 46.0\% in the reverse
        direction.
  \item \textbf{NN initialization contribution}: HNSGA-III vs.\
        NK-HNSGA-III shows nearly identical average cost ($<0.01\%$
        difference) but substantially better stability (Std: 96.5 vs.\
        129.6), HV (0.7995 vs.\ 0.7540), and GD (0.1956 vs.\ 0.2211).
  \item \textbf{Satellite configuration}: The optimal number of active
        satellites varies strongly by day (Sat=1: 43\%, Sat=2: 20\%,
        Sat=3: 37\%), demonstrating the necessity of the multi-period
        formulation.
\end{itemize}

\subsection{Limitations}

(1)~Soft time windows; hard-window extension would require additional
feasibility repair mechanisms.  (2)~Synthetic benchmarks only; real-world
validation remains an important future step.  (3)~Sequential FE-SE
decomposition may yield suboptimal solutions relative to joint optimization
in some cases.  (4)~Homogeneous fleet assumption; heterogeneous or electric
vehicles are not modeled.  (5)~Sample size $n=10$ limits statistical power.

\subsection{Future Research Directions}

(1)~Hard time-window extension with feasibility-aware decoding.
(2)~Real-world validation on empirical logistics datasets.
(3)~Joint FE-SE co-evolutionary or bi-level optimization.
(4)~Green vehicle integration (electric fleets, range constraints, charging
schedules).
(5)~Dynamic/stochastic demand via robust or stochastic optimization.
(6)~Enhanced VNS neighborhoods (e.g., inter-route exchange, cross-satellite
relocation).

% ============================================================
\printbibliography
% ============================================================

\end{document}
